\textbf{Proof.}
Write
\[
 \psi(\vartheta)=\frac{p_+-p_-}{n},
 \qquad
 N_\vartheta=\binom{n}{p_+,p_-,r_0}
 =\frac{n!}{p_+!p_-!r_0!}.
\]
The three counts are nonnegative integers with sum \(n\), so \(N_\vartheta\) is finite and positive, including when any count is zero. By the design requirement, \(A\sim\mathcal D\) is independent of \(\vartheta\), the ordered partition, and the fair bits \((H_i)_{i\in R}\); dependence among the coordinates of \(A\) remains unrestricted.

Fix \(\vartheta\) and \(A=a\). If \(i\in P\), its response type is \((1,0)\), and hence \(S_i=1\) under either assigned arm. If \(i\in N\), its type is \((0,1)\), and hence \(S_i=0\). If \(i\in R\), then the transformation either leaves the fair bit \(H_i\) unchanged or replaces it by \(1-H_i\). Thus the variables \((S_i)_{i\in R}\) are independent fair bits conditional on \((\vartheta,P,N,R,A=a)\), and
\[
 X=p_++\operatorname{Binomial}(r_0,1/2)
\]
conditional on \(\vartheta\), with a law independent of \(a\).

For the stronger kernel assertion, fix \(s\in\{0,1\}^n\), put \(x=\sum_i s_i\), and suppose \(p_+\leq x\leq p_++r_0\). A partition can generate \(S=s\) only if \(P\) is a \(p_+\)-element subset of \(\{i:s_i=1\}\) and \(N\) is a \(p_-\)-element subset of \(\{i:s_i=0\}\); after those choices, \(R\) is forced. There are
\[
 \binom{x}{p_+}\binom{n-x}{p_-}
\]
such ordered partitions, and for each one exactly one realization of its \(r_0\) fair bits produces \(s\). Therefore
\[
 \Pr(S=s\mid\vartheta,A=a)
 =\frac{\binom{x}{p_+}\binom{n-x}{p_-}}
 {N_\vartheta2^{r_0}}.
\]
The right side depends on \(s\) only through \(x\) and not on \(a\). The factorial identity
\[
 \binom nx\binom{x}{p_+}\binom{n-x}{p_-}
 =N_\vartheta\binom{r_0}{x-p_+}
\]
then gives the displayed binomial law and
\[
 \Pr(S=s\mid X=x,\vartheta,A=a)
 =\binom nx^{-1}\mathbf 1\!\left\{\sum_i s_i=x\right\}.
\]
Values outside the support may be defined arbitrarily. The factorization also proves that \(A\) is ancillary and independent of \((S,X,\vartheta)\) under every mixture \(\nu\).

Given \(A\), the maps between \(Y^{\mathrm{obs}}\) and \(S\) are inverse. Thus a full-data estimator is a function \(T(A,S)\). Define
\[
 T_1(s)=\mathbb E_\mathcal D\{T(A,s)\},
 \qquad
 T_2(x)=\binom nx^{-1}\sum_{s:\,\sum_i s_i=x}T_1(s).
\]
Independence of \(A\), the parameter-free conditional law of \(S\) given \(X\), and two applications of conditional Jensen yield
\[
 \mathbb E_\nu\{T_2(X)-\psi(\vartheta)\}^2
 \leq \mathbb E_\nu\{T(A,S)-\psi(\vartheta)\}^2.
\]
Conversely every scalar rule \(f(X)\) is a legitimate full-data rule under every design, because \(X\) is observable from \((A,Y^{\mathrm{obs}})\). Infimizing the two sides proves equality of the full-data and scalar Bayes risks. All factorial and conditional-probability conventions used above remain valid at the stated boundary counts. \(\square\)